% Generated draft: requires mathematical and source review.
\documentclass[11pt]{article}
\usepackage[margin=25mm]{geometry}
\usepackage{fontspec}
\setmainfont{Latin Modern Roman}
\usepackage{amsmath,amssymb,mathtools,unicode-math}
\setmathfont{Latin Modern Math}
\usepackage{enumitem,needspace,xurl,hyperref}
\setlength{\parindent}{0pt}
\setlength{\parskip}{5pt}
\setlength{\emergencystretch}{3em}
\newcommand{\N}{\mathbb{N}}
\newcommand{\Z}{\mathbb{Z}}
\newcommand{\Q}{\mathbb{Q}}
\newcommand{\R}{\mathbb{R}}
\newcommand{\C}{\mathbb{C}}
\newcommand{\F}{\mathbb{F}}
\newcommand{\A}{\mathbb{A}}
\newcommand{\PP}{\mathbb{P}}
\newcommand{\E}{\mathbb{E}}
\newcommand{\eps}{\varepsilon}
\newcommand{\dd}{\,\mathrm{d}}
\begin{document}
\section*{UnsolvedMath ID: 30001590}
\textbf{Source problem code:} catalogue entry
\subsection{Definitions and mathematical statement}
The supplied problem concerns solutions near a family of solitons. It does not specify an evolution equation, a spatial domain, boundary conditions, a solution space, or a particular soliton family. Thus the following notation only organizes the question; it does not supply missing hypotheses.

If the intended equation admits a parameterized soliton family, write its profiles schematically as
\[
\mathcal S=\{S_a:a\in A\},
\]
where the parameter set \(A\) and the meaning of each parameter remain unspecified. Whenever an additive decomposition is meaningful in the intended solution space, write a solution near this family schematically as
\[
u(t)=S_{a(t)}+r(t).
\]
Here \(a(t)\) represents the soliton parameters and \(r(t)\) represents the remaining component, called the ``wave part'' in the source. The source gives no decomposition rule, uniqueness condition, orthogonality condition, or criterion ensuring that this remainder is a wave or radiation component. Its phrase ``modal parameters'' is retained without identifying it with modulation parameters or spectral mode amplitudes.

\textbf{Statement of the supplied research question.} What can be said about solutions near solitons? Describe the interaction between the soliton's modal parameters and the wave part of the solution, including blow-up along the soliton family.

This is a research question requiring further specification, rather than a quantified conjecture for a defined class of equations. In particular, the source does not define the topology or size of the neighborhood of \(\mathcal S\), the admissible initial data, the time interval or asymptotic regime, or whether the desired conclusions concern all nearby solutions or a restricted class. It also does not specify whether ``blow-up along the soliton family'' means concentration of soliton profiles, degeneration of their parameters, divergence of a solution norm, or another mechanism. No finite-time or infinite-time blow-up condition can therefore be imposed from the supplied evidence. The title and supplied literature assessment use soliton-resolution terminology, but the record supplies no precise resolution limit or asymptotic decomposition to assert.
\subsection{Short English statement}
The question asks how solutions close to a soliton family evolve, how the soliton's parameters interact with the wave component, and how blow-up may occur along that family. The equation, notion of closeness, parameter decomposition, and meaning of blow-up are unspecified, so a precise quantified conjecture cannot yet be extracted. This is a statement-only draft requiring human mathematical review.
\subsection{Sources}
\begin{itemize}
\item Local UnsolvedMath registry: permanent ID \(30001590\), source problem code \texttt{catalogue entry}, title ``Soliton Resolution and Blow-Up near Soliton Families,'' in the Oberwolfach Reports Open Problems collection. Repository registry locations are \texttt{data/unsolvedmath/problems.json} and \texttt{catalog/problems.json}; the supplied complete record is the evidence used here. No individual proposer is named. The recorded proposed year is 2010.
\item Supplied report citation: \emph{Nonlinear Waves and Dispersive Equations} (2011), \emph{Oberwolfach Reports}, pp.\ 2393--2463, DOI \texttt{10.4171/owr/2010/41}, \url{https://doi.org/10.4171/owr/2010/41}. The background describes Oberwolfach Reports as documenting mathematical workshops and their open-problem discussions. The supplied report title/year and the record's 2010 year information are preserved without reconciling their dating.
\item Supplied literature source, labeled \texttt{doi.org/10.4171/owr/2010/41}: \url{https://doi.org/10.4171/owr/2010/41}. The literature assessment dated exactly \texttt{2026-08-21} labels the problem open and describes the extracted prompt as a broad soliton-resolution question rather than a settled theorem, with work treating particular equations and regimes only. This is supplied dated curation derived from the OWR record, not an independently verified literature assessment. Its time-sensitive status and cited claims require independent verification.
\item The registry labels the statement ``exact'' and reports that manual review found the extracted text self-contained, mathematically coherent, and free of LaTeX or report-layout artifacts. These are supplied verification claims; the present draft does not independently verify the report text or its mathematical completeness.
\end{itemize}
\end{document}
